%% ch05 --- One line of algebra, every kind of behaviour
%%
%% Written September 2026. The book's hinge: the chapter in which the reader
%% SEES chaos appear from one line of arithmetic, x -> r*x*(1 - x), as the
%% knob r is turned from 2.8 to 3.9. Every number the reader cannot do in
%% their head comes from code/measure/ch05.py, which imports the functions
%% the listings print (code/ch05/), so the page and the listings cannot
%% disagree. The map uses only multiplication and subtraction, so even the
%% values taken from a chaotic orbit are identical on every machine.
%%
%% WHAT THIS CHAPTER MAY NOT SAY. That nearby starts separate is shown once,
%% in a teaser figure with no number on it: counting the steps is Chapter 6's
%% measurement, the rate is Chapter 7's Lyapunov exponent, and the whole
%% picture of r against behaviour (period doubling, Feigenbaum's constant,
%% where the windows sit and what period three implies) is Chapter 8's. The
%% period-3 window at r = 3.83 is shown as a surprise and handed on.
%%
%% Twin: chapters/pl/ch05-logistic-map.tex. Section for section, macro for
%% macro, number for number; tools/parity.py compares them.

\chapter{One line of algebra, every kind of behaviour}
\label{ch:logistic}
\index{logistic map}\index{chaos}

\noindent
Here is a rule you could teach a child. Take a number between $0$ and $1$.
Multiply it by one minus itself, then multiply the answer by a fixed number
$r$. That is the next number. Do the same to it, and again, and again. No
dice, no measurement, no outside world: each step is fixed exactly by the
one before.

Make a guess before you read on. If you turn the rule's one knob, $r$,
slowly upwards, how many different kinds of behaviour do you expect: one,
two, a handful? This chapter turns the knob and watches. By the end you will
have seen a list of numbers that never settles and never repeats, looks as
if it came off a roulette wheel, and was produced by nothing but that line.

\begin{outcomes}
\outcome{iterate the logistic map by hand for a step or two, and in Python
  for as many steps as you like}
\outcome{find the map's steady state for a given $r$, and say from the slope
  there whether it holds}
\outcome{draw a cobweb diagram on a curved rule and read from it what an
  orbit will do}
\outcome{tell a steady state, a cycle and a never-repeating orbit apart, and
  say what a program that looks for cycles can and cannot conclude}
\outcome{say what \emph{deterministic} and \emph{chaotic} mean, and why
  looking random is not the same as being random}
\end{outcomes}

\section{The rule and its one knob}
\label{sec:ch05-rule}
\index{logistic map!rule}

\noindent
Chapter~\ref{ch:feedback} ended with a population whose growth slows near
its limit: at a gentle rate it settled, and at a faster one it overshot and
swung. This chapter writes a rule of that kind in the form that made it
famous. Let $x$ be the population as a fraction of the largest it can be, a
number between $0$ and $1$. Write $x_n$, read \enquote{$x$ sub $n$}, for its
value after $n$ steps, so that $x_0$ is the start. The rule is
\[ x_{n+1} = r\, x_n\, (1 - x_n). \]
In words: the next value is $r$ times the present value times one minus the
present value. The factor $x_n$ says that more parents make more young; the
factor $1 - x_n$ is the feedback, the room left, which shrinks as the
population grows. The number $r$ is the growth knob. A rule that takes a
number to the next number is called a \emph{map}, and this one is the
\emph{logistic map}. Beside the rules of Chapter~\ref{ch:iteration},
$x_{n+1} = a\,x_n + b$, it has one extra multiplication, by $1 - x_n$, and
everything in this chapter comes from that one factor.

The product $x\,(1 - x)$ is largest at $x = \tfrac12$, where it is
$\tfrac12 \times \tfrac12 = \tfrac14$, so the next value is never more than
$r/4$. For $r$ up to $4$, then, $x$ never leaves the range from $0$ to $1$.

\begin{think}
Take $r = \num{2.8}$ and start at $x_0 = \num{0.5}$. What is $x_1$? And
$x_2$? The first you can do in your head; the second needs a pencil at most.
\end{think}
\reveal{$x_1 = \num{2.8} \times \num{0.5} \times \num{0.5} = \num{0.7}$, and
$x_2 = \num{2.8} \times \num{0.7} \times \num{0.3} = \num{0.588}$.}

\noindent
That is all there is to it, and a computer does millions of steps a second.
Here is the rule as the book's code writes it.

\pyregion{code/ch05/four_rates.py}{rule}{The logistic map: one line of arithmetic}{lst:ch05-rule}

\noindent
Multiplication and subtraction are rounded in exactly the same way by every
computer, so the numbers in this chapter come out the same to the last digit
on any machine. For a rule whose whole interest, as you are about to see,
lies in its digits, that matters.

\section{Turning the knob}
\label{sec:ch05-knob}
\index{logistic map!four behaviours}

\noindent
Start every run at $x_0 = \num{0.2}$. Let the rule take $200$ steps without
looking, so that the \emph{orbit} \dash{} the list of values a rule produces
from a start \dash{} has time to settle into whatever it will do, then write
down the next eight values. Do that at four settings of the knob:
$r = \num{2.8}$, $\num{3.2}$, $\num{3.5}$ and $\num{3.9}$.

\begin{think}
Chapter~\ref{ch:feedback} found that a gentle growth rate settles and a
faster one overshoots and swings. These four settings run from moderate to
fast. After $200$ steps, what do you expect at $r = \num{3.9}$: one value
repeated, a swing between a few values, or something else?
\end{think}
\reveal{Something else. At $r = \num{3.9}$ the values never settle, and
never fall into a pattern.}

\pyregion{code/ch05/four_rates.py}{run}{Two hundred steps without looking, then eight more, at four settings}{lst:ch05-run}

\transcript{ch05-four-rates}

\noindent
Each row is a different kind of behaviour. At $r = \num{2.8}$ the population
has come to rest at \val{ch05.fp.28}: a \emph{steady state}. At
$r = \num{3.2}$ it alternates for ever between \val{ch05.two.lo} and
\val{ch05.two.hi}, high year, low year: a \emph{cycle of period two}. At
$r = \num{3.5}$ it goes round four values, from \val{ch05.four.lo} to
\val{ch05.four.hi}, and round the same four again: a cycle of period four.
At $r = \num{3.9}$ the numbers jump about with no pattern you can see, and,
as you will check below, none turns up however long you wait.

\plotfig{ch05-four-series}{The four settings over the first sixty steps,
all from $x_0 = \num{0.2}$. Look at how each panel ends: resting,
alternating, going round four values, and still jumping.}{four time series
of the logistic map}

\noindent
In Figure~\ref{fig:ch05-four-series} every panel starts with a scramble,
and three of the four soon find their rhythm. The fourth never does.

\begin{trapbox}
\textbf{\enquote{A simple rule gives simple behaviour.}} All four panels
come from the same line of arithmetic; only $r$ changed, from $\num{2.8}$
to $\num{3.9}$. The behaviour went from one resting value to numbers that
never repeat. How complicated a system's behaviour looks tells you very
little about how complicated its rule is.
\end{trapbox}

\section{The steady state, and when it lets go}
\label{sec:ch05-fixed}
\index{fixed point}\index{fixed point!stability}

\noindent
A \emph{fixed point} is a value the rule leaves unchanged, and a steady state
is a fixed point the orbit has settled on. The logistic map always has one at
$x = 0$: no population stays no population. For the other, ask when
$r\,x\,(1 - x)$ equals $x$. Divide both sides by $x$ to get
$r\,(1 - x) = 1$, so $1 - x = 1/r$, and
\[ x^{*} = 1 - \frac{1}{r}. \]
The star is the usual mark for a fixed point.

\begin{think}
For $r = 2$ the formula gives $x^{*} = \tfrac12$. Check that the rule really
leaves $\tfrac12$ unchanged. What does the formula give for $r = 4$, and
does that check out too?
\end{think}
\reveal{$2 \times \tfrac12 \times \tfrac12 = \tfrac12$. For $r = 4$ it gives
$\tfrac34$, and $4 \times \tfrac34 \times \tfrac14 = \tfrac34$.}

\noindent
At $r = \num{2.8}$ the formula gives \val{ch05.fp.28}, exactly where the
printout came to rest. At $r = \num{3.2}$ it gives \val{ch05.fp.32}
\dash{} and the orbit swings round it instead, from \val{ch05.two.lo} to
\val{ch05.two.hi}. The fixed point exists; the orbit will not stay on it.

Chapter~\ref{ch:feedback} explained why with a measurement. Nudge the state
a tiny amount off the fixed point and take one step: the nudge comes back
multiplied by the \emph{slope} of the rule there. Between $-1$ and $1$,
nudges shrink and the fixed point holds; beyond, they grow and the orbit is
pushed away. A negative slope also flips the nudge to the other side at
every step. The library's \api{slope} measures it by nudging $x$ a millionth
each way.

\pyregion{code/ch05/fixed_point.py}{fixed}{The fixed point, and the slope of the rule there}{lst:ch05-fixed}

\transcript{ch05-fixed-point}

\noindent
Each slope is $2 - r$: $2 - \num{2.8} = \num{-0.8}$,
$2 - \num{3.2} = \num{-1.2}$, and so on. (With calculus, the slope of
$r\,x\,(1 - x)$ is $r\,(1 - 2x)$, and at $x = 1 - 1/r$ that is exactly
$2 - r$.) At $r = \num{2.8}$ a nudge flips and shrinks at every step, and
the orbit closes in. At $\num{3.2}$ it flips and grows, and the orbit is
thrown off.

\begin{think}
As you turn $r$ up from $\num{2.8}$, at what value does the slope $2 - r$
reach $-1$, so that the fixed point stops holding?
\end{think}
\reveal{At $r = 3$, where $2 - 3 = -1$. Below $3$ the steady state holds;
above $3$ it lets go.}

\noindent
That is the first place where the knob changes the kind of behaviour, and
you found it by hand. The change from two values to four comes further up;
Chapter~\ref{ch:bifurcations} finds where, and a pattern in where.

\begin{lab}{l05_01_stability}{Where the steady state lets go}
Write \code{fixed\_point}; \code{slope\_at}, which measures the slope at a
point by nudging it each way; and \code{holds}, which says whether a nudge
at the fixed point shrinks. Then write \code{lets\_go}, which finds where
\code{holds} turns false between $\num{2.5}$ and $\num{3.5}$ by halving the
interval again and again. It should agree with the answer above.
\end{lab}

\section{The cobweb on a curve}
\label{sec:ch05-cobweb}
\index{cobweb diagram}

\noindent
Chapter~\ref{ch:iteration} drew iteration as a cobweb on a straight line.
The same drawing works for any rule. Draw the rule's graph, the height
$r\,x\,(1 - x)$ for every $x$ from $0$ to $1$: a hump, highest at
$x = \tfrac12$. Draw the diagonal, where the height equals $x$. Start at
$x_0$ on the bottom edge and go up to the hump: the height you reach is
$x_1$. Go across to the diagonal, which puts you above $x_1$; up or down to
the hump for $x_2$; across again; and so on. The hump crosses the diagonal
at the fixed point, where the next value equals the present one. The
library's \api{cobweb} lists the corners of the path.

\pyregion{code/ch05/cobweb_data.py}{corners}{The corners of a cobweb, three steps long}{lst:ch05-cobweb}

\transcript{ch05-cobweb}

\noindent
Each \code{up} lands on the hump at the next value of the orbit,
$\num{2.8} \times \num{0.2} \times \num{0.8} = \num{0.448}$ for the first;
each \code{across} lands on the diagonal at the same height.

\begin{think}
At $r = \num{2.8}$ the cobweb spirals in to the crossing. At $r = \num{3.2}$
the orbit ends up alternating between two values for ever. What shape does
its cobweb settle into?
\end{think}
\reveal{A closed rectangle with two corners on the hump: up to one value,
across, down to the other, across, and back to the start.}

\plotfig{ch05-four-cobwebs}{Cobwebs for the four settings, each on its own
hump with the diagonal. At $r = \num{2.8}$ the path spirals in; at
$\num{3.2}$ and $\num{3.5}$ it starts right beside the crossing, is pushed
away and settles on a closed loop; at $\num{3.9}$ it never
closes.}{four cobwebs on the logistic map}

\noindent
Look at where the hump crosses the diagonal in
Figure~\ref{fig:ch05-four-cobwebs}. At $r = \num{2.8}$ it falls
through the crossing less steeply than the diagonal climbs, and the spiral
tightens. At $\num{3.2}$ it falls more steeply, and a path started beside
the crossing spirals outwards until it is caught on the rectangle. A curve
falling exactly as steeply as the diagonal climbs has slope $-1$: the
drawing and the measurement agree. At $\num{3.5}$ the loop has four corners on the hump. At
$\num{3.9}$ the path never closes, and slowly fills a whole band.

\section{Counting the cycle, and failing to}
\label{sec:ch05-period}
\index{period}\index{cycle}

\noindent
For thousands of numbers, let the computer read the cycle. The library's
\api{period} looks at the last $128$ values of an orbit for the smallest
whole number $p$ such that every value equals the one $p$ steps later, to
within a billionth. It returns that $p$, or Python's word for nothing,
\code{None}, when no $p$ up to $64$ fits.

\pyregion{code/ch05/periods.py}{detect}{Asking for the period after two thousand steps}{lst:ch05-periods}

\transcript{ch05-periods}

\noindent
The detector agrees with the printout: $1$, then $2$ and $4$. At
$r = \num{3.9}$, and at $r = 4$, the top of the range, it finds nothing.

\begin{think}
At $r = \num{3.9}$ the detector found no cycle up to $64$ steps long at the
end of two thousand steps. Suppose you ran the rule for a hundred thousand
steps instead. Would a repeat turn up eventually?
\end{think}
\reveal{No. In a hundred thousand steps not a single value comes back.}

\pyregion{code/ch05/never_repeats.py}{count}{A hundred thousand steps, looking for any value seen before}{lst:ch05-never}

\transcript{ch05-never-repeats}

\noindent
A Python \code{set} keeps one copy of each different value, so a value that
came back would make the second count smaller. It does not: all
\val{ch05.never.values} values are different. That settles the question for
these \val{ch05.never.steps} steps, because if a value ever did come back,
the one after it would too \dash{} the same number always produces the same
next number \dash{} and so would everything after that: a cycle. No value
returning means no cycle has begun. The second line of each pair matters
later: the same start, run again, gives the same numbers exactly.

\begin{trapbox}
\textbf{\enquote{If it has not repeated after a thousand steps, it is about
to.}} A long wait makes a repeat feel overdue, as a coin that has shown heads
five times seems to owe you tails. The orbit owes you nothing. No cycle is
pulling it in, and running longer only produces more numbers that have
never appeared before.
\end{trapbox}

\begin{rigourbox}
A run is evidence, not proof. For $r = 4$ it can be proved that from almost
every start the orbit never repeats; for $r = \num{3.9}$ the evidence is
computation like this. And a computer holds only finitely many numbers
between $0$ and $1$, so a computed orbit must eventually meet a value it has
had before, and cycle \dash{} after more steps than this run took.
Chapter~\ref{ch:computers} says what that means for a simulation.
\end{rigourbox}

\begin{lab}{l05_02_period}{A cycle detector of your own}
Write \code{settle}, which runs the rule without looking and then keeps the
next few values, and \code{period}, which finds the smallest $p$ for which
every value in a list equals the one $p$ places on, and returns \code{None}
when no $p$ up to a limit fits. The tests use made-up lists and this
section's orbits.
\end{lab}

\begin{think}
The detector found a cycle of four at $r = \num{3.5}$ and nothing at
$\num{3.9}$. What do you expect it to find at $r = \num{3.83}$, a little
below $\num{3.9}$?
\end{think}
\reveal{A cycle of period three. After a short scramble the orbit goes round
\val{ch05.win.a}, \val{ch05.win.b} and \val{ch05.win.c} for ever.}
\index{periodic window}

\plotfig{ch05-window}{Order inside the jumping. At $r = \num{3.83}$ the
orbit scrambles for a while, then locks onto a cycle of three; at
$r = \num{3.9}$, close by, it never settles.}{a cycle of three inside the
chaos}

\noindent
This surprises almost everyone, and it should. Turning the knob up does not
make the behaviour steadily wilder: among the settings where orbits jump
without pattern there are narrow windows where order comes back. Where they
sit, and what a cycle of three implies, belong to the whole picture that
Chapter~\ref{ch:bifurcations} draws.

\section{Random-looking, and not random}
\label{sec:ch05-may}
\index{deterministic}\index{chaos!definition}

\noindent
Now the names. A rule is \emph{deterministic} when the present fixes the next
state completely, with nothing left to chance, so the same start always gives
the same orbit. You watched the logistic map be so: run twice from
$\num{0.2}$, it gave the same hundred thousand numbers to the last digit. An
orbit is \emph{chaotic} when it comes from a deterministic rule, stays within
bounds, never settles into rest or a cycle \dash{} and, the part not yet
measured, is so sensitive to its start that any difference there, however
small, eventually grows until the two futures have nothing in common.

\plotfig{ch05-two-starts}{Two runs at $r = \num{3.9}$, from $\num{0.2}$ and
from $\num{0.200001}$, a millionth higher. For a while the lines lie on top
of each other; then they part for good.}{two nearby starts at one setting}

\noindent
That last part is Chapter~\ref{ch:butterfly}'s measurement;
Figure~\ref{fig:ch05-two-starts} is only a preview. Together they make the strangest fact in this book so far. The
orbit at $r = \num{3.9}$ looks random, and is not: every number in it was
fixed exactly by the one before. Yet you cannot foretell it far ahead
without knowing its start more precisely than anyone knows anything.

\index{May, Robert}
In 1976 Robert May, a physicist who had moved into ecology, published a
review in \emph{Nature} whose title says what this chapter has shown:
\emph{Simple mathematical models with very complicated dynamics}. He wrote
it for biologists.

\begin{think}
An ecologist counts the insects in a field every summer for twenty years.
The counts jump up and down with no pattern anyone can find. What does that
tell you about the cause?
\end{think}
\reveal{Less than it seems. The weather might do it, or errors in counting
\dash{} or nothing but the population's own rule of growth.}

\noindent
That third possibility was May's point. Irregular numbers from the field
were usually read as the mark of a noisy environment or of imperfect
counting; the logistic map shows that one generation a year and one simple
rule of growth can produce the same irregularity with no disturbance at all.
May ended with a plea that reaches beyond ecology: that students should
meet this equation early, because an intuition trained only on straight-line
rules leaves people unprepared for what curved ones do. You have now met it.

\begin{trapbox}
\textbf{\enquote{Irregular output needs a random or complicated cause.}}
The $r = \num{3.9}$ row in Section~\ref{sec:ch05-knob} came from the same
line of arithmetic as the rows above it, with no randomness anywhere.
Irregular output can come from chance, from a complicated mechanism, or from
a simple deterministic rule, and the output alone does not say which.
Chapter~\ref{ch:life} teaches how to tell them apart.
\end{trapbox}

\begin{worldbox}
Fisheries science meets the same behaviour in a different formula. William
Ricker's stock-and-recruitment model of 1954, which relates the fish
spawning in one year to the young they produce, is a one-line rule of the
same kind: growth that slows as numbers rise, a generation at a time. Turn
up its growth rate and it passes, in the same order as the logistic map,
from a steady stock through alternating booms and busts to irregular swings.
A catch that swings from year to year does not, by itself, prove that the
weather did it.
\end{worldbox}

\noindent
One line, one knob: a steady state, cycles of two and four, orbits that
never repeat, and islands of order among them. What is still missing is a
measure of how unpredictable the chaos is. Chapter~\ref{ch:butterfly}
counts the steps until two nearby starts part, Chapter~\ref{ch:lyapunov}
turns that into one number, and Chapter~\ref{ch:bifurcations} draws every
setting of the knob in one picture.

\begin{summarybox}
\begin{itemize}
\item The \textbf{logistic map}, $x_{n+1} = r\,x_n\,(1 - x_n)$, is a
  straight-line rule with one extra multiplication, the feedback $1 - x_n$.
  For $r$ up to $4$, $x$ stays between $0$ and $1$.
\item Turning $r$ gives a \textbf{steady state} at $\num{2.8}$, a
  \textbf{cycle of period two} at $\num{3.2}$, one of period four at
  $\num{3.5}$, and at $\num{3.9}$ an orbit that never settles or repeats.
\item The \textbf{fixed point} is $x^{*} = 1 - 1/r$ and the \textbf{slope}
  there is $2 - r$, so the steady state holds for $r$ below $3$.
\item A \textbf{cobweb} spirals in to a fixed point that holds, is pushed
  off one that does not, settles on a closed loop for a cycle, and at
  $\num{3.9}$ never closes.
\item A \textbf{cycle detector} reports $1$, $2$, $4$, then \code{None},
  which means not found, never proved chaotic. No value comes back in a
  hundred thousand steps; at $\num{3.83}$ a \textbf{window} holds a cycle of
  three.
\item \textbf{Deterministic}: same start, same orbit, every time.
  \textbf{Chaotic}: deterministic, bounded, never settling, sensitive to the
  start. It looks random and is not \dash{} May's 1976 warning.
\end{itemize}
\end{summarybox}

\canyou

\begin{checkpoint}
\item With $r = 3$ and $x_0 = \num{0.5}$, what are $x_1$ and $x_2$?
  \answerto{$x_1 = 3 \times \num{0.5} \times \num{0.5} = \num{0.75}$ and
  $x_2 = 3 \times \num{0.75} \times \num{0.25} = \num{0.5625}$.}
\item What is the fixed point other than zero when $r = \num{2.5}$, and does
  it hold?
  \answerto{$1 - 1/\num{2.5} = \num{0.6}$. The slope there is
  $2 - \num{2.5} = \num{-0.5}$, between $-1$ and $1$, so it holds.}
\item At $r = \num{3.2}$ the fixed point is \val{ch05.fp.32}, yet the orbit
  does not stay there. Why not?
  \answerto{The slope there is $2 - \num{3.2} = \num{-1.2}$: a nudge flips
  side and grows at every step, so the orbit is pushed off, onto a cycle of
  two.}
\item A cobweb settles on a closed loop with four corners on the hump. What
  is the orbit doing?
  \answerto{Going round a cycle of period four; the corners are the four
  values it visits in turn.}
\item A cycle detector returns \code{None}. What can you conclude, and what
  can you not?
  \answerto{That no cycle up to the length it checked was found in the
  stretch it looked at. Not that the orbit is chaotic: a longer cycle, or an
  orbit still settling, gives \code{None} too.}
\item Yearly counts of an animal population jump about with no pattern.
  Name three possible causes.
  \answerto{The weather or other chance; errors in counting; and the
  population's own deterministic rule of growth, as in the logistic map at
  $r = \num{3.9}$.}
\item What does it mean that the logistic map is deterministic, and how can
  its output still look random?
  \answerto{The same start always gives exactly the same orbit. At
  $r = \num{3.9}$ it never settles or repeats, so it shows no pattern,
  although every value is fixed by the one before.}
\end{checkpoint}
